\documentclass[11pt]{article}
\usepackage[margin=1in]{geometry}
\usepackage{amsmath,amssymb,amsthm}
\usepackage{algorithm}
\usepackage{algpseudocode}
\usepackage{booktabs}
\usepackage{hyperref}

\newtheorem{definition}{Definition}
\newtheorem{theorem}{Theorem}
\newtheorem{lemma}[theorem]{Lemma}
\newtheorem{property}{Property}

\title{\textbf{Mathematical Specification, Surrogate Jacobian Relaxations, and Bias Analysis of the Straight-Through Estimator (ALGO-NN-29)}}
\author{Team Scaibu \\ \small Email: \texttt{jaydeep.9664920749@gmail.com} \\ \small Architecture and Core Engine Team}
\date{\today}

\begin{document}
\maketitle

\begin{abstract}
Discrete step functions, sign binarizations, and integer roundings have zero derivatives almost everywhere and Dirac delta singularities at jump boundaries, which prevents standard gradient backpropagation. The Straight-Through Estimator (STE) applies the exact discrete operator during the forward pass while substituting a continuous surrogate identity or clipped Jacobian $\frac{\partial q}{\partial x} \approx \mathbb{I}_{|x| \le c}$ during the backward pass. We formalize surrogate gradient mechanics, prove bounded bias properties, provide full pseudocode matching the reference implementation, and derive complexity.
\end{abstract}

\section{Notation and Mathematical Preliminaries}
\label{sec:notation}

Let $x \in \mathbb{R}$ denote a continuous latent variable and $q = Q(x) \in \mathcal{S}_{\mathrm{discrete}}$ its quantized representation (e.g. $\{-1, +1\}$ or $\mathbb{Z}$). Let $\frac{\partial \mathcal{L}}{\partial q}$ denote the upstream error gradient.

\subsection{Surrogate Gradient Mechanics}
\paragraph{Intuitive Meaning:} During training, neural network weights need continuous nudges to learn. STE lets discrete operations (like rounding or thresholding) happen in the forward pass so the model behaves like a quantized system, but pretends the operation was smooth during backpropagation so gradients can pass through.

\paragraph{Formal Mathematical Definition:}
\begin{itemize}
    \item \textbf{Forward Pass:} $q = \operatorname{sgn}(x) = \begin{cases} +1 & \text{if } x \ge 0, \\ -1 & \text{if } x < 0. \end{cases}$
    \item \textbf{Clipped Identity Backward (HardTanh STE):}
    \begin{equation}
    \label{eq:ste_backward}
    \frac{\partial \mathcal{L}}{\partial x} = \frac{\partial \mathcal{L}}{\partial q} \cdot \mathbb{I}_{\{|x| \le c\}}.
    \end{equation}
\end{itemize}

\paragraph{Worked Numerical Example:}
Let $x = 0.5$, upstream gradient $\frac{\partial \mathcal{L}}{\partial q} = 2.0$, clip threshold $c = 1.0$:
Forward: $q = \operatorname{sgn}(0.5) = +1.0$.
Backward: Since $|0.5| \le 1.0$, $\frac{\partial \mathcal{L}}{\partial x} = 2.0 \times 1.0 = 2.0$.
If $x = 1.5$: Forward $q = +1.0$, Backward: $|1.5| > 1.0 \implies \frac{\partial \mathcal{L}}{\partial x} = 2.0 \times 0.0 = 0.0$.

\subsection{Summary of Symbols}
\begin{itemize}
    \item $N \in \mathbb{N}$: Number of elements in the tensor.
    \item $\mathbf{x} \in \mathbb{R}^N$: Continuous latent tensor.
    \item $\mathbf{dy} \in \mathbb{R}^N$: Upstream incoming gradients $\frac{\partial \mathcal{L}}{\partial \mathbf{q}}$.
    \item $c \in \mathbb{R}_{>0}$: Gradient clipping boundary ($c = 1.0$).
    \item $\mathbf{q} \in \mathbb{R}^N$: Quantized discrete forward output.
    \item $\mathbf{dx} \in \mathbb{R}^N$: Surrogate backward gradient $\frac{\partial \mathcal{L}}{\partial \mathbf{x}}$.
\end{itemize}

\section{Problem Definition and Core Concepts}
\label{sec:problem_definition}

\subsection{Quantization Modes}
\begin{itemize}
    \item \textbf{Sign Mode:} $q = \operatorname{sgn}(x) \in \{-1, +1\}$ for 1-bit binary neural networks.
    \item \textbf{Round Mode:} $q = \lfloor x + 0.5 \rfloor \in \mathbb{Z}$ for integer quantization (INT8/INT4).
    \item \textbf{Clamp STE Mode:} $q = \operatorname{clamp}(x, -c, c)$ for bounded activations.
\end{itemize}

\section{Algorithmic Specification}
\label{sec:algorithm}

Algorithm~\ref{alg:ste} specifies the forward-backward STE evaluation matching \texttt{impl.py}.

\begin{algorithm}[H]
\caption{Straight-Through Estimator (STE) Evaluation}
\label{alg:ste}
\begin{algorithmic}[1]
\Require Continuous inputs $\mathbf{x} \in \mathbb{R}^N$, upstream gradients $\mathbf{dy} \in \mathbb{R}^N$, mode $\mathrm{mode} \in \{\text{sign}, \text{round}, \text{clamp\_ste}\}$, threshold $c > 0$.
\Ensure Quantized outputs $\mathbf{q} \in \mathbb{R}^N$, surrogate gradients $\mathbf{dx} \in \mathbb{R}^N$.
\State Initialize $\mathbf{q} \gets [], \mathbf{dx} \gets []$
\For{$i = 1$ to $N$}
    \State $x \gets x_i, \quad dy \gets dy_i$
    \If{$\mathrm{mode} = \text{sign}$}
        \State $q_{\mathrm{val}} \gets 1.0$ if $x \ge 0.0$ else $-1.0$
    \ElsIf{$\mathrm{mode} = \text{round}$}
        \State $q_{\mathrm{val}} \gets \operatorname{round}(x)$
    \ElsIf{$\mathrm{mode} = \text{clamp\_ste}$}
        \State $q_{\mathrm{val}} \gets \max(-c, \min(c, x))$
    \EndIf
    \If{$|x| \le c$}
        \State $dx_{\mathrm{val}} \gets dy$
    \Else
        \State $dx_{\mathrm{val}} \gets 0.0$
    \EndIf
    \State Append $q_{\mathrm{val}}$ to $\mathbf{q}$, append $dx_{\mathrm{val}}$ to $\mathbf{dx}$
\EndFor
\State \Return $\mathbf{q}, \quad \mathbf{dx}$
\end{algorithmic}
\end{algorithm}

\section{Theoretical Guarantees and Bias Analysis}
\label{sec:guarantees}

\begin{theorem}[Surrogate Gradient Energy Preservation]
\label{thm:ste_energy}
\textbf{Assumptions:} Let $x \sim \mathcal{N}(0, \sigma^2)$ be standard normal input and $c \ge 1$.
\textbf{Guarantee:} The expected surrogate gradient norm $\mathbb{E}[\|\mathbf{dx}\|_2^2]$ is strictly non-zero and bounded by $\mathbb{E}[\|\mathbf{dx}\|_2^2] = (1 - 2\Phi(-c/\sigma)) \mathbb{E}[\|\mathbf{dy}\|_2^2] > 0$, preventing gradient collapse through discrete layers.
\textbf{Proof:}
Since $\mathbf{dx} = \mathbf{dy} \cdot \mathbb{I}_{\{|x| \le c\}}$, taking expectations under independence between upstream error and layer input:
\begin{equation}
\mathbb{E}[\|\mathbf{dx}\|_2^2] = \mathbb{E}[\|\mathbf{dy}\|_2^2] \cdot P(|x| \le c) = \mathbb{E}[\|\mathbf{dy}\|_2^2] \left( \Phi\left(\frac{c}{\sigma}\right) - \Phi\left(-\frac{c}{\sigma}\right) \right) = \mathbb{E}[\|\mathbf{dy}\|_2^2] \left( 1 - 2\Phi\left(-\frac{c}{\sigma}\right) \right).
\end{equation}
Since $c > 0, \sigma > 0 \implies \Phi(-c/\sigma) < 0.5 \implies 1 - 2\Phi(-c/\sigma) > 0$.
\textbf{Limitations:} STE introduces a biased gradient estimator ($\mathbb{E}[\hat{\mathbf{g}}] \neq \nabla f$); optimization trajectories can oscillate around quantization boundaries.
\end{theorem}

\section{Complexity Analysis}
\label{sec:complexity}

\begin{itemize}
    \item \textbf{Time Complexity:} $\mathcal{O}(N)$ comparison and assignment operations.
    \item \textbf{Space Complexity:} $\mathcal{O}(N)$ memory for quantized and gradient buffers.
\end{itemize}

\section{Worked Numerical Example}
\label{sec:worked_example}

Let $\mathbf{x} = [-1.5, 0.2, 0.8], \mathbf{dy} = [1.0, 1.0, 1.0], c = 1.0$ in sign mode:
\begin{enumerate}
    \item $x_1 = -1.5$: $q_1 = -1.0$. Since $|-1.5| > 1.0 \implies dx_1 = 0.0$.
    \item $x_2 = 0.2$: $q_2 = +1.0$. Since $|0.2| \le 1.0 \implies dx_2 = 1.0$.
    \item $x_3 = 0.8$: $q_3 = +1.0$. Since $|0.8| \le 1.0 \implies dx_3 = 1.0$.
    \item Outputs: $\mathbf{q} = [-1.0, 1.0, 1.0]$, $\mathbf{dx} = [0.0, 1.0, 1.0]$.
\end{enumerate}

\section{Numerical Considerations and Edge Cases}
\label{sec:numerical}

\begin{itemize}
    \item \textbf{Symmetric Threshold:} Gradient clipping at $|x| \le c$ prevents weights from growing unbounded outside the active quantization dynamic range.
\end{itemize}

\begin{thebibliography}{9}
\bibitem{bengio2013} Y.~Bengio, N.~L{\'e}onard, and C.~Courville, ``Estimating or Propagating Gradients Through Stochastic Neurons for Conditional Computation,'' \emph{arXiv:1308.3432}, 2013.
\end{thebibliography}

\end{document}
